% ============================================================
% Section 4: Information Budget (Medium length for TNNLS)
% Target: ~2.5 pages
% ============================================================

\section{The Information Budget}
\label{sec:budget}

We now formalize an empirical diagnostic framework for understanding when GNNs can improve over feature-only methods.

\subsection{Core Definitions}

\begin{definition}[Empirical Enhancement Ceiling]
\label{def:budget}
Given a node classification task with features $X$ and labels $Y$, under a fixed experimental protocol, we define the \textbf{Empirical Enhancement Ceiling} as:
\begin{equation}
    \mathcal{B} = 1 - \text{Acc}_{\text{MLP}}^*(X, Y)
\end{equation}
where $\text{Acc}_{\text{MLP}}^*$ is the best accuracy achievable using only node features.
\end{definition}

\textbf{Intuition}: The MLP accuracy represents the predictive power of features alone. The ceiling $\mathcal{B}$ quantifies the \textit{maximum room for improvement} that graph structure could potentially provide.

\subsection{The Conditional Enhancement Bound}

\begin{theorem}[Conditional Enhancement Bound]
\label{thm:conditional_bound}
Let $\mathcal{P}$ denote a fixed experimental protocol. If the following conditions hold:
\begin{enumerate}
    \item[(C1)] \textbf{Fair Comparison}: MLP and GNN trained with equivalent capacity
    \item[(C2)] \textbf{Consistent Features}: Both use identical input features $X$
    \item[(C3)] \textbf{Standard Aggregation}: GNN uses neighborhood aggregation without external signals
\end{enumerate}
Then the GNN improvement over MLP satisfies:
\begin{equation}
    \text{Acc}_{\text{GNN}} - \text{Acc}_{\text{MLP}} \leq \mathcal{B} = 1 - \text{Acc}_{\text{MLP}}
\end{equation}
\end{theorem}

\begin{proof}
Rearranging the claimed inequality gives $\text{Acc}_{\text{GNN}}\leq 1$, which holds by definition of accuracy. Conditions (C1)--(C3) are experimental comparability requirements, not mathematical requirements for this arithmetic ceiling.
\end{proof}

\subsection{Empirical Content Beyond the Arithmetic Bound}

The ceiling $\Delta \leq 1-\text{Acc}_{\text{MLP}}$ follows arithmetically from $\text{Acc}_{\text{GNN}} \leq 1$. Its scientific content therefore lies not in observing zero violations, but in testing whether headroom supports calibrated model-selection predictions.

\textbf{(1) Shared-Term Audit.} On the nine traceable external datasets, the correlation between $\mathcal{B}$ and best-candidate-GNN advantage is $r=-0.255$ ($R^2=0.065$; 95\% bootstrap CI $[-0.743,0.523]$). A null model that permutes GNN accuracies while retaining the shared $-\text{Acc}_{\text{MLP}}$ term produces mean $r=0.669$ (95\% interval $[0.327,0.897]$). Thus the previously reported large positive correlation is not reproduced by the audited external result set and is not used as evidence in this revision.

\begin{table}[t]
\centering
\caption{Shared-Term Audit on Nine External Datasets}
\label{tab:predictive_comparison}
\begin{tabular}{@{}lcc@{}}
\toprule
Quantity & Estimate & 95\% interval \\
\midrule
Observed Pearson $r$ & $-0.255$ & $[-0.743,0.523]$ \\
Shared-term null $r$ & $0.669$ (mean) & $[0.327,0.897]$ \\
\bottomrule
\end{tabular}
\end{table}

\textbf{Interpretation} (Table~\ref{tab:predictive_comparison}): The arithmetic ceiling remains valid, but the audited data do not support using raw Budget--advantage correlation as independent predictive evidence. Model selection must instead be evaluated directly on held-out datasets with a fixed prediction target.

\textbf{(2) LODO Cross-Validation.} Leave-one-dataset-out validation gives 12/12 correct predictions among decisive cases (95\% CI: 73.5\%--100\%) at 75\% coverage (12/16); four cases are explicitly treated as abstentions.

\textbf{(3) Structural Intervention.} Degree-preserving edge shuffle experiments show that disrupting the observed topology changes GNN advantage from $+12.6\%$ to $-34.8\%$ on Cora (47.4 point drop).

\textbf{(4) Non-Vacuous Predictions.} Unlike vacuous bounds (e.g., ``$\text{Acc} \leq 1$''), the Budget enables actionable predictions: when $\mathcal{B} < 5\%$, all 7 high-MLP datasets show $|\Delta| < 5\%$ (Prediction 4 in Table~\ref{tab:predictions}).

\subsection{Falsifiable Predictions}

\begin{table}[t]
\centering
\caption{Four Falsifiable Predictions and Their Validation}
\label{tab:predictions}
\small
\begin{tabular}{@{}p{0.4\linewidth}p{0.5\linewidth}@{}}
\toprule
Prediction & Validation \\
\midrule
P1: $\Delta \leq \mathcal{B}$ (arithmetic ceiling) & Audited subset satisfies \\
P2: $\Delta$ correlates with $\mathcal{B}$ & Not supported on audited external set ($r=-0.255$) \\
P3: Edge shuffle decreases $\Delta$ & Cora: $+12.6\% \to -34.8\%$ \\
P4: $\mathcal{B} < 5\% \Rightarrow |\Delta| < 5\%$ & All 7 high-MLP datasets satisfy \\
\bottomrule
\end{tabular}
\end{table}

\subsection{Ceiling-Efficiency Decomposition}

The Information Budget enables a decomposition of GNN improvement:
\begin{equation}
    \text{GNN Advantage} = \underbrace{(1 - \text{Acc}_{\text{MLP}})}_{\text{Ceiling } \mathcal{B}} \times \underbrace{\eta(h, \kappa)}_{\text{Efficiency } \in [-1, 1]}
\end{equation}
where efficiency $\eta$ depends on graph properties (homophily $h$, SNR ratio $\kappa$).

\begin{table}[t]
\centering
\caption{Efficiency Varies with Homophily (Cora, Fixed Budget)}
\label{tab:efficiency}
\begin{tabular}{@{}ccccc@{}}
\toprule
$h$ & MLP & Budget & GCN Adv. & Efficiency \\
\midrule
0.81 (original) & 75.1\% & 24.9\% & +12.7\% & \textbf{51.0\%} \\
0.55 (rewired) & 75.1\% & 24.9\% & +3.0\% & 12.0\% \\
0.35 (rewired) & 75.1\% & 24.9\% & $-9.2\%$ & $-36.9\%$ \\
\bottomrule
\end{tabular}
\end{table}

\textbf{Key Insight} (Table~\ref{tab:efficiency}): Same budget, vastly different outcomes. Homophily determines efficiency; features determine the ceiling.

\subsection{Practical Implications}

\begin{enumerate}
    \item \textbf{Always train MLP first}: If $\text{Acc}_{\text{MLP}} > 95\%$, the budget is too small ($< 5\%$) for meaningful GNN improvement.

    \item \textbf{Predict GNN potential}: High budget ($\mathcal{B} > 20\%$) combined with high homophily ($h > 0.7$) suggests significant GNN benefit.

    \item \textbf{Diagnose failures}: When GNN underperforms MLP despite high budget, investigate graph quality (low $h$, noisy edges).
\end{enumerate}

\subsection{Information-Theoretic Justification}

\begin{proposition}[Structure Information Bound]
\label{prop:structure_bound}
Let $P_e^{\text{MLP}}$ and $P_e^{\text{GNN}}$ denote the Bayes-optimal error rates. Then:
\begin{equation}
    P_e^{\text{MLP}} - P_e^{\text{GNN}} \leq \frac{I(Y; G | X)}{\log C}
\end{equation}
where $I(Y; G | X)$ is the conditional mutual information between labels and structure given features.
\end{proposition}

This formalizes why GNN improvement is bounded by information in the graph not already captured by features.
